\section{An exact decomposition into scales}\label{sec:layers}
The large-order proof uses the scale decomposition of~\cite{OAI2026}, with
explicit parameters, a parity correction, and three final reduction rates.
We specify which edges and which vertex multiplicities are counted.

A supported piece consists of a graph on the original vertex set together
with a vertex support containing both endpoints of every one of its edges.
Its mass is the size of that support. Different pieces may share vertices;
their edge sets must be disjoint. The mass of a family is the sum of its
support sizes, including multiplicity.

\begin{lemma}[One scale]\label{lem:scale}
At logarithmic scale $w$, set
\[
 D=2^w,\quad \ell=2^{(1+\ep)w},\quad L_0=2^{\rho w},\quad h_0=2^{sw}.
\]
Under the parameter inequalities at the applicable scale, a piece with
average degree at most $D$ can be partitioned into long cycles, even
expanding cores, and a new family of supported remainder pieces. The
following bounds hold per unit parent mass:
\begin{enumerate}[label=(\roman*),leftmargin=*]
\item the number of long cycles is at most $D/(2\ell)$;
\item the sum of box orders is at most
      $1+2\ell/(L_0-\ell-2)$, and each box has order at most $L_0$;
\item cores in a box have disjoint vertex supports, orders greater than
      $h_0$, and cut expansion at least $(1023/1024)h_0$; every core degree
      is even;
\item the next average-degree bound is at most
      $h_0(1001/1000+\log_2(L_0/h_0))$.
\end{enumerate}
\end{lemma}
\begin{proof}
First remove actual edge-disjoint cycles longer than the prescribed
threshold; their lengths pay (i). On the cycle-free remainder, a DFS
partition with integer parameters $\floor\ell$ and
$\floor{(L_0-\floor\ell)/2}$ gives the boxes in (ii). In particular, the
denominator includes both rounding units.
Within each box, repeatedly cut a side of insufficient edge expansion.
The entropy potential pays the cut edges and gives the logarithmic term
in (iv). The surviving induced cores have disjoint supports. For each core,
choose a forest whose odd-degree vertices are exactly the odd-degree
vertices of the core, and move that forest into the remainder. Subtraction
makes the core even. The cut bound for the forest and the certified scale
inequalities leave the expansion in (iii); its at most one edge per vertex
is absorbed by the $1001/1000$ term in (iv).
The remainder children are defined \emph{after} this move. Each operation
is a partition of original edges, so composition preserves the partition.
The conditions $\ell+2<L_0$, $h_0>4096$, $h_0\le L_0$, and
$2048L_0<h_0^2$ are checked on the complete scale ranges, not assumed from
asymptotic notation.
\end{proof}

Let $X=\log_2n\ge w_A$. Choose $k\ge0$ so that
\begin{equation}\label{eq:phase}
           w_A\le W:=\theta_2^kX<w_A/\theta_2,
 \qquad m_0=k+3.
\end{equation}
(The symbol $m_0$ counts stages, not edges.) Define
\begin{equation}\label{eq:schedule}
 w_i=\begin{cases}
       \theta_2^iX,&0\le i\le k+1,\\
       \theta_1\theta_2^{k+1}X,&i=k+2,\\
       \theta_0\theta_1\theta_2^{k+1}X,&i=k+3.
      \end{cases}
\end{equation}
Apply Lemma~\ref{lem:scale} at stages $0,\ldots,k+2$. The entropy
inequalities for the three transitions give average degree at most
$2^{w_{i+1}}$ in every child piece. The last processed scale is at least
$q_0w_A=w_0$. Let $S_i$ be the total core mass at stage $i$.

For each prefix, the long cycles removed so far, the cores extracted so far,
and the \emph{whole} current frontier form an exact partition of $E(G)$.
This prefix identity permits both operations used later: replacing one
selected core layer together with its descendant frontier, or stopping
before the last two stages. It would not suffice to retain only the sparse
edges at the end of all stages.

The overlap factor in (ii) is at most $1+\chi2^{-\delta w}$.
Products of these factors are bounded by
$\prod_j(1+x_j)\le(1-\sum_jx_j)^{-1}$ when $x_j\ge0$ and
$\sum_jx_j<1$. The certified endpoint inequalities give a uniform mass
factor $M$. For a phase lower endpoint $a\le W$, define
\begin{align}
 M_1(a)&=\frac1{1-\chi(2^{-\delta q_0a}+E(\delta\theta_2a))},\notag\\
 M_2(a)&=\frac1{1-\chi E(\delta\theta_2a)},&
 M_3(a)&=\frac1{1-\chi E(\delta a)}.\label{eq:phase-masses}
\end{align}
All denominators are positive. In particular
\begin{equation}\label{eq:last-core-masses}
 S_{m_0-1}\le M_1(a)n,\qquad S_{m_0-2}\le M_2(a)n.
\end{equation}

\begin{lemma}[Stopped remainder]\label{lem:stopped}
For $a\le W\le d$, stop the actual partition before stage $k+1=m_0-2$.
Merge its frontier pieces on the original vertex set and call the resulting
graph $H$. Then
\[
     2\abs{E(H)}\le M_3(a)2^{\theta_2d}n.
\]
The long cycles retained before this stop number at most
\begin{equation}\label{eq:Jtrim}
                   J_{\rm tr}(a)n,
       \qquad J_{\rm tr}(a)=\tfrac12M_3(a)E(\ep a).
\end{equation}
\end{lemma}
\begin{proof}
The frontier average-degree bound is $2^{w_{k+1}}=2^{\theta_2W}$
per unit support mass. Its support mass is at most $M_3(a)n$.
The pieces have disjoint edges, so summing their degree bounds proves the
first inequality after the merge. Sum the long-cycle estimates from
Lemma~\ref{lem:scale}(i) over stages $0,\ldots,k$; their scales are bounded
below in reverse order by $a\theta_2^{-j}$. The same mass bound gives
\eqref{eq:Jtrim}.
\end{proof}

Every retained even core also has a completion packing of weight at most
$(1+C\eta)$ times its support size, provided the weighted order bound is
available for strictly smaller graphs. The completion construction uses
an edge bank, actual paths, and a smaller-order remainder. We state the
more precise version, including the extra credit needed for folding, next.
